\documentclass[11pt]{article}
\newcommand{\doctitle}{Response to Collaborator Goals}
\newcommand{\shorttitle}{Trial Analysis Guide}
\newcommand{\docdate}{2026-03-22}
% pmsimstats-ng standard document preamble
% Usage: % pmsimstats-ng standard document preamble
% Usage: % pmsimstats-ng standard document preamble
% Usage: \input{pmsimstats-preamble} after \documentclass[11pt]{article}
%
% Documents must define before \input:
%   \newcommand{\doctitle}{...}       Full document title
%   \newcommand{\shorttitle}{...}     Short title for header
%   \newcommand{\docdate}{YYYY-MM-DD} Document date

\usepackage[margin=1in]{geometry}
\usepackage{amsmath,amssymb}
\usepackage[T1]{fontenc}
\usepackage{lmodern}
\usepackage{microtype}
\usepackage{graphicx}
\graphicspath{{figures/}}
\usepackage{booktabs}
\usepackage{longtable}
\usepackage{array}
\usepackage{hyperref}
\usepackage{xcolor}
\usepackage{fancyhdr}
\usepackage{parskip}
\usepackage{caption}
\usepackage{enumitem}
\usepackage{verbatim}
\usepackage{listings}

\lstset{
  language=R,
  basicstyle=\small\ttfamily,
  keywordstyle=\color{blue!70!black},
  commentstyle=\color{green!50!black},
  stringstyle=\color{red!60!black},
  breaklines=true,
  frame=single,
  framerule=0.5pt,
  rulecolor=\color{gray!40},
  backgroundcolor=\color{gray!5},
  xleftmargin=1em,
  framexleftmargin=1em,
  columns=flexible
}

\hypersetup{
  colorlinks=true,
  linkcolor=blue!60!black,
  citecolor=green!50!black,
  urlcolor=blue!60!black
}

\pagestyle{fancy}
\fancyhf{}
\lhead{\footnotesize pmsimstats-ng Technical Report}
\rhead{\footnotesize \shorttitle}
\rfoot{\footnotesize \thepage}
\renewcommand{\headrulewidth}{0.4pt}

\title{\doctitle}
\author{pmsimstats team}
\date{\docdate}
 after \documentclass[11pt]{article}
%
% Documents must define before \input:
%   \newcommand{\doctitle}{...}       Full document title
%   \newcommand{\shorttitle}{...}     Short title for header
%   \newcommand{\docdate}{YYYY-MM-DD} Document date

\usepackage[margin=1in]{geometry}
\usepackage{amsmath,amssymb}
\usepackage[T1]{fontenc}
\usepackage{lmodern}
\usepackage{microtype}
\usepackage{graphicx}
\graphicspath{{figures/}}
\usepackage{booktabs}
\usepackage{longtable}
\usepackage{array}
\usepackage{hyperref}
\usepackage{xcolor}
\usepackage{fancyhdr}
\usepackage{parskip}
\usepackage{caption}
\usepackage{enumitem}
\usepackage{verbatim}
\usepackage{listings}

\lstset{
  language=R,
  basicstyle=\small\ttfamily,
  keywordstyle=\color{blue!70!black},
  commentstyle=\color{green!50!black},
  stringstyle=\color{red!60!black},
  breaklines=true,
  frame=single,
  framerule=0.5pt,
  rulecolor=\color{gray!40},
  backgroundcolor=\color{gray!5},
  xleftmargin=1em,
  framexleftmargin=1em,
  columns=flexible
}

\hypersetup{
  colorlinks=true,
  linkcolor=blue!60!black,
  citecolor=green!50!black,
  urlcolor=blue!60!black
}

\pagestyle{fancy}
\fancyhf{}
\lhead{\footnotesize pmsimstats-ng Technical Report}
\rhead{\footnotesize \shorttitle}
\rfoot{\footnotesize \thepage}
\renewcommand{\headrulewidth}{0.4pt}

\title{\doctitle}
\author{pmsimstats team}
\date{\docdate}
 after \documentclass[11pt]{article}
%
% Documents must define before \input:
%   \newcommand{\doctitle}{...}       Full document title
%   \newcommand{\shorttitle}{...}     Short title for header
%   \newcommand{\docdate}{YYYY-MM-DD} Document date

\usepackage[margin=1in]{geometry}
\usepackage{amsmath,amssymb}
\usepackage[T1]{fontenc}
\usepackage{lmodern}
\usepackage{microtype}
\usepackage{graphicx}
\graphicspath{{figures/}}
\usepackage{booktabs}
\usepackage{longtable}
\usepackage{array}
\usepackage{hyperref}
\usepackage{xcolor}
\usepackage{fancyhdr}
\usepackage{parskip}
\usepackage{caption}
\usepackage{enumitem}
\usepackage{verbatim}
\usepackage{listings}

\lstset{
  language=R,
  basicstyle=\small\ttfamily,
  keywordstyle=\color{blue!70!black},
  commentstyle=\color{green!50!black},
  stringstyle=\color{red!60!black},
  breaklines=true,
  frame=single,
  framerule=0.5pt,
  rulecolor=\color{gray!40},
  backgroundcolor=\color{gray!5},
  xleftmargin=1em,
  framexleftmargin=1em,
  columns=flexible
}

\hypersetup{
  colorlinks=true,
  linkcolor=blue!60!black,
  citecolor=green!50!black,
  urlcolor=blue!60!black
}

\pagestyle{fancy}
\fancyhf{}
\lhead{\footnotesize pmsimstats-ng Technical Report}
\rhead{\footnotesize \shorttitle}
\rfoot{\footnotesize \thepage}
\renewcommand{\headrulewidth}{0.4pt}

\title{\doctitle}
\author{pmsimstats team}
\date{\docdate}


\begin{document}
\maketitle
\thispagestyle{fancy}
\tableofcontents

Date: 2026-03-23\\
Context: Following the pmsimstats code audit and revision

\section{Code Locations}

All functions have been consolidated into the
\texttt{pmsimstats-ng} repository:

\begin{table}[h]
\centering
\small
\begin{tabular}{lll}
\toprule
Function & File & Purpose \\
\midrule
\texttt{analyze\_trial\_extended()} & \texttt{R/carryover\_analysis.R} & Main drug effect + interaction + predicted curves \\
\texttt{characterize\_carryover()} & \texttt{R/carryover\_analysis.R} & Sweep half-lives, compare AIC, find best fit \\
\texttt{print\_trial\_summary()} & \texttt{R/carryover\_analysis.R} & Clinical summary of trial analysis \\
\texttt{print\_carryover\_summary()} & \texttt{R/carryover\_analysis.R} & Summary of carryover characterization \\
\texttt{lme\_analysis()} & \texttt{R/lme\_analysis.R} & Core LME analysis (nlme + corCAR1) \\
\texttt{generateData()} & \texttt{R/generateData.R} & Data generation with cached sigma \\
\texttt{buildSigma()} & \texttt{R/generateData.R} & Covariance matrix construction \\
\texttt{validateParameterGrid()} & \texttt{R/generateData.R} & PD pre-validation \\
\texttt{generateSimulatedResults()} & \texttt{R/generateSimulatedResults.R} & Full simulation loop with parallel processing \\
\bottomrule
\end{tabular}
\end{table}

\bigskip\noindent\rule{\textwidth}{0.4pt}\bigskip

\section{Goal 1: Did prazosin improve PTSD symptoms?}

\textit{How to analyze this in this trial design? -- RCH implement
with RGT guidance}

\subsection{What we have}

The revised \texttt{lme\_analysis()} uses \texttt{nlme::lme} with \texttt{corCAR1}
(continuous-time AR(1) residual correlation), which is
directly applicable to actual trial data. Vignette 3 in the
original package demonstrates applying \texttt{lme\_analysis()} to
real clinical trial data (\texttt{CTdata.rda}). The model formula
for the N-of-1 design is:

\begin{verbatim}
Sx ~ bm + t + Dbc + bm:Dbc, random = ~1|ptID,
  correlation = corCAR1(form = ~t|ptID)
\end{verbatim}

\subsection{What is needed}

The current model tests the biomarker-treatment
\textit{interaction} (\texttt{bm:Dbc}): does blood pressure predict who
responds? To test whether prazosin improves symptoms
\textit{overall}, the coefficient of interest is \texttt{Dbc} (the main
drug effect), not \texttt{bm:Dbc}. This is a simpler test -- check
whether the drug indicator is significant.

The existing infrastructure supports this; it only requires
extracting a different coefficient from the model output.

\subsection{Action -- DONE}

Implemented in \texttt{R/carryover\_analysis.R} as
\texttt{analyze\_trial\_extended()}. Extracts both the main drug
effect (\texttt{Dbc}) and the biomarker interaction (\texttt{bm:Dbc})
with 95\% CIs, variance components, predicted response
curves, and clinical decision thresholds.

Usage:
\begin{lstlisting}
result <- analyze_trial_extended(td, dat, op, threshold = 5)
print_trial_summary(result)
\end{lstlisting}

\bigskip\noindent\rule{\textwidth}{0.4pt}\bigskip

\section{Goal 2: Characterize carryover in actual data}

\textit{Need to figure out how to structure the carryover
quantification and identify the best sequence of response
measures. -- RCH implement with RGT guidance}

\subsection{What we have}

\begin{itemize}
  \item The \texttt{lambda\_cor = ln(2)/t\_half} derivation provides a
    principled framework for relating carryover half-life to
    the correlation structure (see
    \texttt{pdf/08-biomarker-correlation-decay.pdf}).
  \item The \texttt{Dbc} continuous drug indicator with exponential decay
    is already implemented and parameterized by
    \texttt{carryover\_t1half}.
  \item The \texttt{nlme::lme} with \texttt{corCAR1} analysis model handles the
    AR(1) residual structure.
\end{itemize}

\subsection{What is needed}

\subsubsection{2a. Structure the carryover quantification}

Fit the actual trial data with multiple carryover models
(different half-lives) and compare fit via AIC/BIC or
likelihood ratio test:

\begin{lstlisting}
results <- list()
for (t_half in c(0.25, 0.5, 1.0, 2.0, 4.0)) {
  # Compute Dbc with this half-life
  # Fit nlme::lme with corCAR1
  # Record AIC, BIC, log-likelihood
  results[[as.character(t_half)]] <- list(
    t_half = t_half,
    aic = AIC(fit),
    bic = BIC(fit),
    loglik = logLik(fit)
  )
}
# Select best-fitting half-life
\end{lstlisting}

This is a straightforward extension of the existing
\texttt{lme\_analysis()}: loop over candidate half-lives and
compare model fit. The \texttt{Dbc} computation is already
parameterized by \texttt{carryover\_t1half}.

A new function \texttt{characterize\_carryover()} would:

\begin{enumerate}
  \item Accept the trial data and design specification
  \item Sweep a range of candidate half-lives
  \item Fit the LME model at each half-life
  \item Return a comparison table (half-life, AIC, BIC,
    log-likelihood, Dbc coefficient, p-value)
  \item Identify the best-fitting half-life
\end{enumerate}

\subsubsection{2b. Identify best sequence of response measures}

This likely means: at which timepoints does the carryover
effect manifest most strongly? The existing trial design
structure (with \texttt{tod}, \texttt{tsd}, \texttt{tpb} at each timepoint)
provides the timing information. The analysis would:

\begin{enumerate}
  \item Fit the model with the best carryover half-life
  \item Extract per-timepoint residuals
  \item Examine whether residuals at off-drug timepoints
    systematically differ from expectations (indicating
    carryover signal)
  \item Plot the residual pattern across timepoints to
    visualize where carryover is strongest
\end{enumerate}

\subsection{Action -- DONE}

Implemented in \texttt{R/carryover\_analysis.R} as
\texttt{characterize\_carryover()}. Sweeps candidate half-lives,
fits nlme::lme at each, compares AIC/BIC, and returns the
best-fitting half-life with full model output.

Usage:
\begin{lstlisting}
cc <- characterize_carryover(td, dat,
  half_lives = c(0.25, 0.5, 1.0, 2.0, 4.0))
print_carryover_summary(cc)
# cc$best_t_half gives the optimal half-life
# cc$best_model gives the full model fit
\end{lstlisting}

Validated: correctly recovers the true t\_half = 1.0 from
simulated data generated with that half-life.

\bigskip\noindent\rule{\textwidth}{0.4pt}\bigskip

\section{Goal 3: Analysis methods to add carryover}

\subsection{3a. Simulation power calculations}

\textit{Done, partially; perhaps needs to be more modular/flexible.
-- RGT implement}

\subsubsection{Status: Essentially complete}

The revised simulation code provides:

\begin{itemize}
  \item \texttt{lambda\_cor} parameter for correlation decay
    (auto-derived from half-life or manual override for
    sensitivity analysis)
  \item Carryover half-life as a simulation parameter:
    \texttt{c(0, 0.5, 1.0)} weeks (clinically realistic values
    replacing the publication's \texttt{c(0, 0.1, 0.2)})
  \item \texttt{Dbc} continuous drug indicator with exponential decay
    in the analysis model
  \item 1,000-replicate validated power estimates with proper
    Type I error control (2.5--5.8\% across all designs)
  \item 8-minute runtime (200 reps) enabling rapid sensitivity
    analysis; 47 minutes for 1,000 reps
  \item Zero positive definiteness failures across all 162
    parameter combinations
\end{itemize}

Under Architecture~A
(\texttt{dgp\_architecture = 'mean\_moderation'}), the
$c_{bm}$ ceiling is removed, allowing evaluation of the full
Hendrickson publication parameter grid including
$c_{bm} = 0.6$.

\subsubsection{Remaining modularity/flexibility items}

\begin{enumerate}
  \item \textbf{Carryover decay model options.} Currently exponential
    only. The \texttt{pmsimstats2025} repo has linear and Weibull
    decay in \texttt{calculate\_carryover()}. These could be
    backported to allow sensitivity analysis over decay
    shape.
  \item \textbf{Separate correlation decay rate.} The \texttt{lambda\_cor}
    parameter currently defaults to \texttt{ln(2)/t\_half}
    (matching the drug elimination rate). Making it an
    explicit sensitivity parameter allows exploring
    scenarios where biomarker predictive power decays
    faster or slower than the drug effect.
  \item \textbf{With/without carryover in analysis model.} The
    simulation could compare power under two analysis
    scenarios: (a) carryover modeled in the analysis
    (continuous Dbc), and (b) carryover ignored (binary
    Db). This addresses the question: how much power is
    gained by correctly modeling carryover in the analysis?
  \item \textbf{Modular design specification.} The trial designs are
    currently hardcoded in the simulation scripts. A
    configuration file or function that generates design
    specifications from high-level parameters (number of
    phases, block durations, randomization points) would
    make it easier to explore alternative designs.
\end{enumerate}

\subsection{3b. Actual analysis}

\textit{In progress, reviewed 2/14 how to get the information
needed in the output structure. -- RGT implement}

\subsubsection{What we have}

The \texttt{lme\_analysis()} with \texttt{full\_model\_out = TRUE} returns:

\begin{itemize}
  \item \texttt{form}: the model formula
  \item \texttt{fit}: the full \texttt{nlme::lme} model object
  \item \texttt{data}: the analysis-ready data with all derived
    variables (\texttt{Dbc}, \texttt{tsd}, \texttt{De}, etc.)
  \item \texttt{summary}: the standard output (beta, betaSE, p-value)
\end{itemize}

The full model object provides access to:

\begin{itemize}
  \item \texttt{summary(fit)\$tTable}: all coefficients with SEs and
    p-values
  \item \texttt{fitted(fit)}: predicted values per observation
  \item \texttt{residuals(fit)}: residuals per observation
  \item \texttt{ranef(fit)}: random effects (per-participant intercepts)
  \item \texttt{intervals(fit)}: confidence intervals for all parameters
  \item \texttt{VarCorr(fit)}: variance components (random intercept
    variance, residual variance, AR(1) parameter)
\end{itemize}

\subsubsection{What is needed for the output structure}

\begin{enumerate}
  \item \textbf{Estimated carryover half-life} from the model
    comparison in Goal 2 (the best-fitting \texttt{t\_half}).
  \item \textbf{Predicted response curves per participant} using the
    fitted model: \texttt{predict(fit, newdata = ...)} with
    participant-specific random effects.
  \item \textbf{Biomarker-specific treatment effect as a continuous
    function of blood pressure.} The \texttt{bm:Dbc} coefficient
    gives the slope: for each unit increase in blood
    pressure, the drug effect changes by $\beta_4$ CAPS
    points. Combined with the main \texttt{Dbc} effect, the
    predicted drug effect for a participant with blood
    pressure \texttt{bm} is:
    \texttt{effect(bm) = beta\_Dbc + beta\_bm:Dbc * (bm - mean\_bm)}
  \item \textbf{Clinical decision threshold.} At what blood pressure
    value does the predicted benefit exceed a clinically
    meaningful threshold (e.g., 5-point CAPS improvement)?
    This is: \texttt{bm\_threshold = mean\_bm +
    (threshold - beta\_Dbc) / beta\_bm:Dbc}
\end{enumerate}

\subsection{Action}

\textbf{DONE.} Implemented as \texttt{analyze\_trial\_extended()} in
\texttt{R/carryover\_analysis.R}. Items 1--4 are all included:
estimated carryover half-life (via \texttt{characterize\_carryover}),
predicted response curves (\texttt{result\$predicted\_effect(bm)}),
biomarker-specific treatment effect with CIs, and clinical
decision threshold (\texttt{result\$threshold\_bm}).

\bigskip\noindent\rule{\textwidth}{0.4pt}\bigskip

\section{Summary of Actions}

\begin{table}[h]
\centering
\small
\begin{tabular}{llll}
\toprule
Goal & Status & Action & Location \\
\midrule
1. Prazosin efficacy & \textbf{DONE} & \texttt{analyze\_trial\_extended()} & \texttt{R/carryover\_analysis.R} \\
2a. Carryover quantification & \textbf{DONE} & \texttt{characterize\_carryover()} & \texttt{R/carryover\_analysis.R} \\
2b. Best response sequence & Conceptual & Residual analysis at each timepoint & Needs implementation \\
3a. Simulation power & \textbf{DONE} & Full simulation with optimizations & \texttt{R/generateSimulatedResults.R} \\
3b. Actual analysis output & \textbf{DONE} & Extended output with predictions & \texttt{R/carryover\_analysis.R} \\
\bottomrule
\end{tabular}
\end{table}

\bigskip\noindent\rule{\textwidth}{0.4pt}\bigskip

\section{Key Deliverables from the Audit}

The following documents and code are available:

\subsection{White Papers}

\begin{itemize}
  \item \texttt{pdf/04-revised-power-analysis.pdf} (13 pages): Complete
    revision summary for the authors, with before/after
    Figure 4 heatmaps and all code changes
  \item \texttt{pdf/18-response-parameter-sensitivity.pdf} (6 pages):
    Response parameter sensitivity analysis
  \item \texttt{pdf/08-biomarker-correlation-decay.pdf} (10 pages):
    Pharmacokinetic derivation of the correlation decay rule
  \item \texttt{pdf/06-ar1-residual-correlation.pdf} (8 pages): Analysis
    model transition from lmer to nlme with corCAR1
  \item \texttt{pdf/07-positive-definiteness-failures.pdf} (10 pages):
    Positive definiteness failure analysis
  \item \texttt{13-figure4-walkthrough.md}: Complete technical walkthrough
    of Figure 4 methodology and code
\end{itemize}

\subsection{Simulation Scripts}

\begin{itemize}
  \item \texttt{analysis/figure4/04-generate-power-head.R}: Figure 4 simulation (current code)
  \item \texttt{analysis/figure4/01-generate-power-sweep.R}: Figure 4 power sweep (sequential)
  \item \texttt{analysis/figure4/02-generate-power-sweep-original.R}: Publication reproduction
  \item \texttt{analysis/figure4/03-generate-power-by-commit.R}: Per-commit comparison
  \item \texttt{analysis/figure4/05-pd-diagnostics.R}: PD failure instrumentation
  \item \texttt{analysis/figure4/06-plot-power-heatmaps.R}: Heatmap plotting
  \item \texttt{analysis/figure5/01-generate-parameter-sensitivity.R}: Figure 5 simulation (revised)
  \item \texttt{analysis/figure5/02-generate-parameter-sensitivity-original.R}: Figure 5 (original)
\end{itemize}

\subsection{Key Results}

\begin{itemize}
  \item 1,000-replicate Figure 4 with zero PD failures
  \item Type I error: 2.5--5.8\% (all nominal)
  \item N-of-1 carryover sensitivity confirmed at realistic
    half-lives (0.5--1.0 weeks)
  \item Design ranking preserved: N-of-1 > CO > OL+BDC > OL
  \item Runtime: 8 minutes (200 reps), 47 minutes (1,000 reps)
\end{itemize}

\vfill
\noindent\rule{\textwidth}{0.4pt}
{\footnotesize
Rendered on 2026-03-25 at 19:01 PDT.\\
Source: \textasciitilde/prj/alz/10-pmsimstats-ng/pmsimstats-ng/docs/14-trial-analysis-implementation-guide.tex}
\end{document}
